\section{Validation with Real Data (GuIA/Sebrae Platform)}\label{sec:real}

The synthetic families of Section~\ref{sec:experiments} embody a hypothesis: that e-learning data are correlated, with longer modules carrying more content. This section tests that hypothesis, and the behavior of the algorithms, on anonymized data from GuIA, an adaptive learning platform in production operated by Sebrae. The experiments are implemented in \code{experiment\_real.py}, which reads the anonymized files in \code{real\_data/} through \code{real\_instances.py}; reproducing them does not require access to the platform's infrastructure.

\subsection{Extraction and anonymization}

The data were extracted from the platform's production database through read-only, fully aggregated queries (a temporary job, destroyed after use). No personal data or text left the database: courses and modules were converted into indices, and the versioned files contain only counts, medians and means. The aggregated source comprises 4 courses, 19 modules, 86 topics, 3,864 enrollments, 576 module sessions with recorded start and end, 189 NPS scores and 13,573 graded exercises.

Two instances were built from it:
\begin{itemize}[leftmargin=1.4em,itemsep=2pt]
\item \textbf{Modules ($n=12$).} The time of a module is the real median observed time, from 16.5 to 376 minutes, a $23\times$ variation that the platform's estimated workload (a constant 2~h per module) does not capture. Three score definitions are used: the completion rate (\emph{conclusao}), the number of topics (\emph{conteudo}) and the course NPS (\emph{nps}). Only the 12 of the 19 modules that have a measured time enter the instance, and times are rounded to the nearest whole minute (minimum 1), because the DP requires integer times.
\item \textbf{Topics ($n=78$).} The time of a module is split among its topics in proportion to the exercise volume of each topic, an observed proxy that is a declared modeling choice, like the score definitions. Two scores are used: the mean exercise grade (\emph{dominio}) and the number of attempts (\emph{engajamento}). The 78 of the 86 topics that have a module time and a recorded mean grade enter the instance, with the same rounding.
\end{itemize}

\subsection{Position of the real data among the synthetic families}

We measured the Pearson correlation between time and score ($r=+1$: they grow together; $r=0$: no relation; $r=-1$: inverse relation). Table~\ref{tab:pearson} shows the result, together with the $p$-value and the 95\% confidence interval computed with SciPy (Section~\ref{sec:validation}). For reference, the synthetic families of Section~\ref{sec:method} have a mean correlation of $+0.11$ (uncorrelated), $+0.98$ (weakly correlated) and $+1.00$ (strongly correlated) over 30 instances.

\begin{table}[h]
\centering
\caption{Pearson correlation between time and score in the real instances.}\label{tab:pearson}
\small
\begin{tabularx}{\linewidth}{@{}l r r r l X@{}}
\toprule
Instance / score & $n$ & $r$ & $p$ & 95\% CI & Reading\\
\midrule
modules / content & 12 & $+0.92$ & $<0.001$ & $[+0.73,+0.98]$ & confirms the hypothesis: a longer module is worth more\\
modules / completion & 12 & $-0.62$ & 0.03 & $[-0.88,-0.08]$ & a longer module is completed less; a regime that no synthetic family covers\\
modules / NPS & 12 & $-0.05$ & 0.89 & $[-0.60,+0.54]$ & no relation\\
topics / mastery & 78 & $+0.13$ & 0.25 & $[-0.09,+0.34]$ & almost no relation\\
topics / engagement & 78 & $+0.29$ & 0.01 & $[+0.07,+0.48]$ & weak relation\\
\bottomrule
\end{tabularx}
\end{table}

Two observations follow. First, the hypothesis of the article that e-learning data are correlated is confirmed for the content-based score ($r=+0.92$), which lies close to the weakly correlated family and below the strongly correlated one. Second, the completion score produces a \emph{negative} correlation, a regime that the three synthetic families do not cover. With $n=12$ the estimates are uncertain, as the confidence intervals show.

\subsection{Agreement of the exact algorithms and greedy loss with short budgets}

On the module instance we ran all algorithms for 3 scores and 4 weekly budgets (2, 5, 10 and 20~h): the exact algorithms found the same optimal value in all 12 combinations. Table~\ref{tab:real-greedy} shows the optimum and the loss of the greedy. A finer sweep of the budget (30~min to 24~h in steps of 15~min) shows that the greedy loses up to 21.0\% with a 30-minute budget and 17.3\% with 1~h (completion score), and 18.8\% with 1.5~h (content score); with 20~h the loss is zero for all three scores. At the topic level ($n=78$) the worst greedy gap drops to 7.3\% (engagement score, $T=45$~min), which confirms on real data the prediction that the gap decreases with $n$ (Section~\ref{sec:exp-greedy}).

\begin{table}[h]
\centering
\caption{Real module instance ($n=12$): optimal score and greedy loss for the weekly budgets.}\label{tab:real-greedy}
\small
\begin{tabular}{@{}l rr rr rr rr@{}}
\toprule
 & \multicolumn{2}{c}{$T=2$~h} & \multicolumn{2}{c}{$T=5$~h} & \multicolumn{2}{c}{$T=10$~h} & \multicolumn{2}{c}{$T=20$~h}\\
\cmidrule(lr){2-3}\cmidrule(lr){4-5}\cmidrule(lr){6-7}\cmidrule(lr){8-9}
Score & opt. & loss & opt. & loss & opt. & loss & opt. & loss\\
\midrule
completion & 238 & 0.0\% & 420 & 4.5\% & 565 & 0.0\% & 619 & 0.0\%\\
content & 21 & 9.5\% & 37 & 5.4\% & 56 & 3.6\% & 79 & 0.0\%\\
NPS & 386 & 0.0\% & 671 & 0.4\% & 953 & 0.0\% & 1143 & 0.0\%\\
\bottomrule
\end{tabular}
\end{table}

\subsection{Granularity: cost and fidelity}

We repeated the pseudo-polynomial test of Section~\ref{sec:exp-pseudo} with the real instance (content score, 10~h budget), changing the time unit. Table~\ref{tab:granularity} shows that only the DP is affected by the unit: it costs 132 operations in hours and 7,212 in minutes, while the other algorithms vary little and without a trend (the range of Branch and Bound is 124 to 348 operations), because changing the unit alters the rounded times but not the size of the problem. Even so, 7,212 operations correspond to milliseconds.

\begin{table}[h]
\centering
\caption{Real module instance (content score, 10~h budget): effect of the time unit.}\label{tab:granularity}
\small
\begin{tabular}{@{}l r r r r r r r@{}}
\toprule
Unit & $T$ & DP ops & DFS ops & B\&B ops & MITM ops & Optimum & Greedy\\
\midrule
hours & 10 & 132 & 4,437 & 124 & 420 & 49 & 49\\
15-min blocks & 40 & 492 & 5,034 & 348 & 408 & 57 & 54\\
5-min blocks & 120 & 1,452 & 4,935 & 311 & 411 & 56 & 54\\
minutes & 600 & 7,212 & 4,884 & 311 & 414 & 56 & 54\\
\bottomrule
\end{tabular}
\end{table}

The test also revealed a loss of fidelity. Rounding the times to hours changed the optimal answer itself: 49 points in hours against 56 in minutes. In hours, the selected path occupies only 502 of the 600 real minutes; in 15-minute blocks the value rises to 57, but the path adds up to 606 minutes and exceeds the budget. A coarse unit therefore brings no relevant saving and distorts the problem; the DP with times in minutes is both cheap and more faithful.

\subsection{One run for all budgets}

The last row of the DP table delivers the optimum for \emph{any} budget from 0 to 24~h at once (Section~\ref{sec:design-dp}). On the real data a single table gives, for the completion score, optimal values of 238, 420 and 565 for 2, 5 and 10~h, and all modules fit from 1,108 minutes on. In a product this answers ``what is the best path with 2, 5 or 10~h per week?'' with one run instead of three.

\subsection{Topic-level instance ($n=78$)}

At $n=78$ the DFS and Meet-in-the-Middle are infeasible by design ($2^{79}$ nodes and $2^{39}$ subsets), which is exactly the regime where only DP, Branch and Bound and the heuristics apply. Table~\ref{tab:topics} shows that DP costs exactly $n(T+1)$ operations (9,438 to 93,678 as $T$ grows from 2 to 20~h), while Branch and Bound is cheaper but varies with the score definition (568 to 6,602 operations), and the greedy loses at most 0.9\% at the weekly budgets.

\begin{table}[h]
\centering
\caption{Real topic instance ($n=78$): optimum, greedy loss and operations.}\label{tab:topics}
\small
\begin{tabular}{@{}l r rrr rrr@{}}
\toprule
 & & \multicolumn{3}{c}{mastery} & \multicolumn{3}{c}{engagement}\\
\cmidrule(lr){3-5}\cmidrule(lr){6-8}
Budget & DP ops & opt. & greedy loss & B\&B ops & opt. & greedy loss & B\&B ops\\
\midrule
2~h & 9,438 & 2,659 & 0.1\% & 1,806 & 9,061 & 0.9\% & 568\\
5~h & 23,478 & 3,196 & 0.0\% & 2,358 & 11,710 & 0.0\% & 1,177\\
10~h & 46,878 & 3,611 & 0.0\% & 3,775 & 12,977 & 0.0\% & 2,700\\
20~h & 93,678 & 3,905 & 0.0\% & 6,602 & 13,573 & 0.0\% & 6,567\\
\bottomrule
\end{tabular}
\end{table}
